\cotefile{157-3}
\scanpage{98}
\mspage{47}
à une \deleted{\unsure{fibri}} Cat. \unsure{fibrée|fibré} de la forme $\deleted{\unsure{X_0}}Z\times S\overset{\mathrm{pr}_S}{\unsure{\relbar|\longrightarrow}}S$, \inserted{avec $Z$ dans Cat.}
\deleted{\unsure{Il reste} à interpréter l'inclusion \[X_s\longrightarrow X\qquad\text{dans }\mathrm{Hot}_W\] en termes de l'élément $\unsure{\xi|g|\zeta}$ de $\mathrm{HOT}_W^{\mathrm{lc}}(S)$ défini par $X$.}
\deleted{\unsure{Or}} pour $X_1=Z\times S\longrightarrow S$, il est immédiat que $X_{1s}\longrightarrow X_1$ est dans $W$. Car \deleted{\unsure{$Z$}} un produit cartésien de flèches $\in W$ (à savoir $\mathrm{id}_Z$ et $s\longrightarrow S$) est dans $W$. (En fait, l'inclusion est $=$ \unsure{un|une} homotopisme !)
Notons maintenant que, si $X_0,X_1$ sont deux \deleted{que l'isomorphie de $X$.} objets de $\underline{\mathrm{Fib}}_0 S$, \deleted{et si l'} \unsure{tels} qu'il existe une $S$-flèche $X_0\longrightarrow X_1$ \inserted{qui soit dans $W_S^d$}, alors $X_{0s}\longrightarrow X_0$ est dans $W$ \unsure{sss|ssi} $X_{1s}\longrightarrow X_1$ l'est. En effet, \unsure{on a}
\diagram{\begin{tikzcd}[ampersand replacement=\&]
X_{0s}\arrow[r,"{\in W}"]\arrow[d,"{i_{\deleted{\unsure{s}}}^{X_0}}"'] \& X_{1s}\arrow[d,"{i_s^{X_1}}"] \\
X_0\arrow[r,"{\in W}"] \& X_1
\end{tikzcd}}
\margin{ici on utilise \unsure{sans doute} $W(2)$, non seulement $W(2\,\mathrm{bis})$,}
qui prouve le résultat.
\deleted{Donc la propriété ``$X_s\longrightarrow X$ est dans $W$'' \inserted{pour un $X$ dans $\underline{\mathrm{Fib}}_0 S$} est stable par la relation d'homotopie \deleted{\unsure{élémentaire},} \inserted{\deleted{dans $\underline{\mathrm{Fib}}_0(S)$ \unsure{ind}} \unsure{rel} $W_S^d$}, (donc si elle est vraie pour un}
Il est vrai que le fait que $X\mathrel{\unsure{\rightleftarrows|=|et}}X_0$ et $X_1=Z\times S$ soient isom. dans $\mathrm{HOT}_W^{\mathrm{lc}}(S)$ n'implique pas, par lui-même, l'existence d'une
% archive sheet number: 98
\scanpage{99}
\mspage{48}
chaîne d'isom. provenant de flèches dans $W_S^d$. Mais en fait, si \deleted{\unsure{on compare}} dans l'énoncé du cor. 2 ci-dessus on compare $X$ dans $\underline{\mathrm{Fib}}_0 S$ à $g^*f^*(X)$, grâce à une homotopie
$\mathrm{id}_S\sim fg$, on trouve bel et bien une suite de foncteurs $\varphi_i:S\longrightarrow S$, avec $\varphi_0=\mathrm{id}_S,\ \varphi_N=fg$, et $\varphi_i$ et $\varphi_{i+1}$ reliés par une flèche \unsure{$\psi$} dans $\underline{\mathrm{Hom}}(S,S)$, donc \deleted{\unsure{$\varphi_i$}} une suite \unsure{des} objets
$\deleted{\unsure{Y}}X_i=\varphi_i^*(X)$ de $\underline{\mathrm{Fib}}_0 S$ (avec $X_0=\unsure{X|x}$, $X_N=g^*f^*(X)$), et $X_i$ et $X_{i+1}$ reliés par une flèche \uline{dans} $\underline{\mathrm{Fib}}_0(S)$ (pas seulement dans $\mathrm{HOT}_W^{\mathrm{lc}}(S)$), laquelle flèche est dans $W_S^d=W_S^f$. Donc on gagne, sous réserve de vérifications de routine diverses. Donc \unsure{modulo|moyennant} celles-ci, on aura
\margin{\unsure{OK}}

\uline{Proposition 9} \deleted{\unsure{(a)}} Supposons que $X$ dans $\mathrm{Cat}_{/S}$ \deleted{soit \unsure{lisse}} \inserted{\deleted{$W$}} \deleted{\unsure{$X'$} \inserted{$\in W$} simplement $W$-prélisse, et que} en soit une $W$-fibration à gauche (cf prop. 8) Alors c'est même une $W$-fibration (au sens de b) page 40,
i.e. \deleted{elle} c'est \unsure{appelé} ``$W$-fibration \unsure{hoto-forte|h.t.-forte}''), en d'autres termes, si $S''\longrightarrow S'$ est \deleted{\unsure{une}} sur $S$, avec $S'',S'$ \unsure{contractibles}, alors $X''\unsure{\relbar|\longrightarrow}X'$ est \deleted{à gauche \unsure{et à} \illegible{1} \unsure{homotopie}} dans $W$. Voire, \inserted{mieux :} si $S''\longrightarrow S'$ est \unsure{un|une} homotopisme, alors $X''\longrightarrow X'$ est dans $W$.
\margin{on trouve seulement que $\mathrm{hot}_W(X'')\longrightarrow\mathrm{hot}_W(X')$ iso.}
% slanted
\margin{C'est OK en supposant que $W$ est \uline{fortement} saturé i.e. $\mathrm{hot}_W(f)$ iso $\Longrightarrow f\in W.$}
% slanted
\margin{\unsure{Rmq} $W$-fibr. à g. $\Longleftrightarrow W$-fibr. à drte $\Longleftrightarrow W$-fibr.}
% archive sheet number: 99
\scanpage{100}
\mspage{49}
Prouvons (si nous pouvons) ce complément \unsure{sympa}. Je précise que j'ai eu à faire l'hypothèse de $W$-prélissité \inserted{sur $X/S$} (ou de $W$-prépropreté), pour être
sûr que l'image inverse $X\times_S S'$, par $f:S'\unsure{\relbar|\longrightarrow}S$, coïncide bien avec $f^*$ pour les $\mathrm{HOT}_W^{\mathrm{lc}}$ -- autrement je n'ai aucune chance. ``On sait'' (tout est à revoir !) que c'est OK
dans les trois cas $X$ $W$-lisse, $X$ $W$-propre, $X$ une $W$-fibration (mais en quel sens ???). \unsure{La} vérification à faire \inserted{(quand on travaille avec les $W_S^d$)} est celle\unsure{-| }ci :
\[
X\overset{i_X}{\longrightarrow}\Phi_S(X)=Y\ \text{est dans }W_S^d\ \text{(sans mérite), et}\tag*{(*)}
\]
le reste par tout changement de base.

Mais pour des $X,Y\in\mathrm{Ob}\,\mathrm{Cat}_{/S}$ qui sont $W$-pré-lisses, les $W_S^d$-équivalences coïncident avec $W_S^f$, laquelle est préservée par chgt de base, ainsi que l'hyp. sur $X,Y$. Donc on est tranquille.

On est ramené à prouver que si $X$ est $W$-prélisse \inserted{sur $S$} \deleted{et si} et une $W$-pseudofibration, \deleted{alors} et si $S'\longrightarrow S$ est un \unsure{homotopisme}, alors $X'\longrightarrow X$ est dans $W$. Par l'argument de \deleted{\unsure{tout}} tantôt utilisant $X\longrightarrow\Phi_S(X)$, on se ramène au cas où $X$ est fibrée scindée
\deleted{Prenant \unsure{une|un} inverse homotopique $g$ de $f$, considérons $g^*(X')=(\underbrace{fg}_{\varphi})^*(X)=X_{\unsure{1|1'}}$}
\deleted{\diagram{\begin{tikzcd}[ampersand replacement=\&]
S'\arrow[r,"{\unsure{f}}"] \& S\arrow[r,"g"] \& S'\arrow[r,"f"] \& S,
\end{tikzcd}}}
% archive sheet number: 100
\scanpage{101}
\mspage{50}
\deleted{et le diagramme multicartésien correspondant}
\deleted{\diagram{\begin{tikzcd}[ampersand replacement=\&]
X'_1\arrow[r]\arrow[d] \& X_1\arrow[r]\arrow[d] \& X'\arrow[r]\arrow[d] \& X=\unsure{/|1}\arrow[d] \\
 S^{\unsure{''|'}}\arrow[r,"f"] \& S\arrow[r,"g"] \& S'\arrow[r,"f"] \& S
\end{tikzcd}}}

La flèche $X'=f^*(X)\unsure{\longrightarrow|\text{---}}X$ s'interprète dans $\mathrm{Hot}_W$ comme
\diagram{\begin{tikzcd}[ampersand replacement=\&, column sep=small]
X'\arrow[rr]\arrow[d,no head] \& \& X\arrow[d,no head] \\
S'\arrow[rr,"f" description,no head]\arrow[dr,"p'"',no head] \& \& S\arrow[dl,"p",no head] \\
\& e \&
\end{tikzcd}}
\[
p'_!(f^*(\unsure{\xi|\zeta}))\longrightarrow p_!(\xi)\qquad\text{où }\xi=\mathrm{hot}_{S,W}(\unsure{X}),\tag*{(*)}
\]
Je dis que cette flèche est iso, i.e. que pour ts $y\in\mathrm{Hot}_W$, la flèche transposée
\[
\deleted{\unsure{\mathrm H}}\mathrm{Hom}_W(p_!(\xi),y)\longrightarrow\mathrm{Hom}(p'_!(f^*(\xi)),y)\tag*{(**)}
\]
est \unsure{bijection|bijective}. Admettant \unsure{la formule|les formules} d'adjonction, la flèche \unsure{(**)} s'interprète comme
\[
\mathrm{Hom}_W(\xi,\underbrace{p^*(y)}_{y_S})\longrightarrow\mathrm{Hom}_W(f^*(\xi),\underbrace{p'{}^*(y)}_{f^*(y_S)})
\]
i.e.
\[
\mathrm{Hom}_W(\xi,y_S)\longrightarrow\mathrm{Hom}_W(f^{\unsure{*|e}}(\xi),f^*(y_{\unsure{S|s}})),
\]
que ce soit iso résulte donc de $f^*$ pl. fid. sur les $\mathrm{HOT}_W^{\mathrm{lc}}$. Or on a vu que $f^*$ est une équivalence (cor. 2. page 46).
% archive sheet number: 101
\scanpage{102}
\mspage{51}
Je vais \unsure{reprendre|refondre} les prop. 9 \inserted{et 8}, un peu mal foutues.

\deleted{\uline{Proposition 9.}} \inserted{Théorème 3.} On suppose que $W$ satisfait à $W(1,2,3)$. Soit $f:X\longrightarrow S$. Considérons les conditions suivantes

a) \deleted{\unsure{$f$ est une $W$}} $\forall\ S''\unsure{\longrightarrow|\text{---}}S'$ sur $S$, avec $S'',S'$ \uline{contractibles}, $X''\longrightarrow X'$ est dans $W$

a$'$) $\forall\ S'$ sur $S$, avec $S'$ contractible, et $s'\in S'$, l'inclusion
\[
X_{s'}\hookrightarrow X'=X\times_S S'
\]
est dans $W$.

$a^{\deleted{\unsure{'|1}}}_1)$ Comme a$'$) mais en supposant que $S'$ a un objet final (au lieu de $S'$ contractible).

$a_2)$ Comme a$'$), mais en supposant que $S'$ a un objet initial.

$b_1)$ $f$ est \inserted{$W$-}\uline{prépropre} (i.e. a$'$ \deleted{\unsure{est sa}} est satisfait quand $s'$ est un objet final de $S'$, sans autre condition sur $S'$ dans $\mathrm{Cat}_{/S}$), et de plus pour toute flèche $s_0\longrightarrow s_1$ dans $S$, la flèche induite $X_{/s_0}\longrightarrow X_{/s_1}$ est dans $W$ (i.e. a$'$) est satisfait pour $Y_{/s_0}\longrightarrow Y_{/s_1}$).

$b_2)$ $f$ est $W$-\uline{prélisse} (i.e. a$'$ est satisfait quand $s'$ est un objet initial de $S'$, sans autre cond. sur $S'$), et de plus les flèches $ {}_{s_0\backslash}X\longrightarrow{}_{s_1\backslash}X$ (pour $s_0\to s_1$ dans $S$) sont dans $W$, i.e. a$'$) satisfaite pour $ {}_{s_1\backslash}Y\longrightarrow{}_{s_0\backslash}Y$).

$c_1)$ $f$ est \inserted{$W$-}prépropre\deleted{\unsure{,}} et \deleted{\unsure{soit}} une $W$-pseudofibration, i.e. la condition a$'$) est satisfaite dans les deux cas suivants : 1°) $s'$ est un objet final de $S'$, 2°) $S'=\Delta1$.

$c_2)$ $f$ est $W$-prélisse et une $W$-pseudofibration, i.e. a$'$) dans les deux cas 1°) $s'$ est objet initial de $S'$, 2°) $S'=\Delta1$.

Toutes les 8 conditions sont équivalentes. Elles impliquent la suivante :

d) Soit $S''\overset{u}{\longrightarrow} S'$ un homotopisme dans $\mathrm{Cat}_{/S}$, \deleted{Alors} et considérons $\deleted{\unsure{X''\longrightarrow}}u_X:X''\longrightarrow X'$. Alors $\mathrm{hot}_W(u_X)$ \deleted{\unsure{est}} est iso,
(Donc si $W$ est fortement \unsure{saturée|saturé}, $u_X\in W$. \deleted{\unsure{E}} Donc sous cette condition \inserted{sur $W$}, d) équivaut aux conditions précédentes.)
% archive sheet number: 102
\scanpage{103}
\mspage{52}
\uline{Définition 4} : On dit que $f:\deleted{\unsure{S}}X\longrightarrow Y$ dans $\mathrm{Cat}$ est une $W$-fibration, si elle satisfait les conditions \inserted{\unsure{équiv.|préc.} a) à $c_2$} de la prop. 9.

\uline{Corollaire 1} Si $f$ est \inserted{$W$-}propre ou \inserted{$W$-}lisse, pour que $f$ soit une $W$-fibration, il faut et \unsure{s|s.} que $f$ soit une $W$-pseudofibration, i.e. que pour $S'=\Delta_1\longrightarrow S$, l\deleted{\unsure{'ori} \unsure{des}}'inclusion des deux fibres
\diagram{\begin{tikzcd}[ampersand replacement=\&]
X'_0\arrow[r,"i_0"] \& X' \& X'_1\arrow[l,"i_1"']
\end{tikzcd}}
de $X'=X\times_S S'$ soient dans $W$.

\uline{Corollaire 2} Supposons $f$ Cat-fibrant. Alors $f$ est une $W$-fibration sss pour \unsure{ts} flèche $u$ \inserted{$:s\longrightarrow s'$} dans $S$, $u^*:X_{s'}\longrightarrow X_s$ est dans $W$. Énoncé dual si $f$ est Cat-cofibrant.

\uline{Corollaire 3} Supposons $f$ $W$-prépropre et $W$-prélisse (par exemple $f$ $W$-parfait). Alors $f$ est une $W$-fibration.
% margin note written sideways
\margin{\uline{Prop.} 10 Soit $f:X\to Y$ une flèche dans \deleted{\unsure{$\mathrm{Cat}$}} $\underline{\mathrm{Fib}}_W S^*$ (sous-cat. pleine de $\mathrm{Cat}_{/S}$ formée des $W$-fibrés sur $S$). Si $\forall s\in S$, $f_s:X_s\to Y_s$ est dans $W$, alors $f\in W$. \deleted{\unsure{Il suffit de}} Considérons \diagram{\begin{tikzcd}[ampersand replacement=\&,column sep=small,row sep=small] X\arrow[r,"f"]\arrow[dr,"{i_X}"']\&Y\arrow[dr,"{i_Y}"]\&{}\\ \&\Phi_S X\arrow[r,"{\Phi_S(f)}"']\&\Phi_S Y\end{tikzcd}}, si $i_X,i_Y,f\in W$, alors $\Phi_S(f)\in W$ \deleted{$\Rightarrow$ \unsure{$f$ est de}} $\Rightarrow f\in W$. Or le fait que $X,$\deleted{\unsure{$Y$}} $W$-fibrés \deleted{\unsure{, sur}} $S$ implique que $i_X\in W_S^f$ (il suffit $=$ que $X$ soit prépropre sur \unsure{$S$}) idem pour $Y,i_Y$. Donc l'hyp. $f\in W_S^f$ implique $\underline{\Phi}_S(f)\in W_S^f$. Donc $\Phi_S(f)\in W$ par prop. 5 (page 29).}
% margin note written sideways
\margin{$^*$ $X$ et $Y$ tous deux prépropres, ou tous deux prélisses \unsure{suffit|suffisent}}

\deleted{\uline{Théorème 4} Soit $W\subset\mathrm{Fl}(\mathrm{Cat})$ satisfaisant $W(1,2,3)$. Considérons les conditions suivantes sur $W$ :}
\deleted{a) Toute $W$-fibration est ultraforte, i.e. \inserted{si $f:X\to S$ dans $\mathrm{Cat}$,} les conditions équiv. de la prop. 9 sur $f$ impliquent que pour $S''\overset{u}{\unsure{\relbar|\longrightarrow}}S'$ dans $\mathrm{Cat}_{/S}$, si $u\in W$, alors $\inserted{u_X:}X''\longrightarrow X'$ est $\unsure{\in}W$.}
\deleted{a$'$) \deleted{Toute \unsure{$W$-fib}} Tout $f$ \inserted{$:X\to S$,} qui est $W_{\unsure{\infty}}$-parfait (donc a fortiori $W_{\deleted{\unsure{\infty}}}$-parfait, donc une $W$-fibration), est une $W$-fibration ultraforte.}
\deleted{b) Pour tte $S$ dans $\mathrm{Cat}$ et toute flèche $f:X\to Y$ dans $\mathrm{Cat}_{/\deleted{\unsure{S}}S}$, si $X,Y$ sont $W$-fibrés sur $S$, et $f\in W$,}
% archive sheet number: 103
\scanpage{104}
\mspage{53}
\uline{IV} \uline{Sur l'axiome} $W(4)$ \uline{et ses variantes}.
% slanted
\margin{Cf aussi \unsure{Partie} \uline{VI} ($\mathrm{HOT}_W$) p. 186–196}

J'ai passé des jours à essayer de prouver que toutes ces variantes sont équivalentes entre elles, sans y avoir réussi. Je vais donner une liste \unsure{de|des} \deleted{\unsure{celles qui}} parmi ces variantes, qui avec le recul me paraissent encore avoir une utilité. (Avec des notations différentes de celles de la p. 13–14)

\tikz[baseline=(n.base)]\node[draw,ellipse,inner sep=2pt](n){$W(4a)$}; Toute $W$-fibration (cf th. 3 p. 51, déf. 3 p. 52) est ultraforte, ce qui revient à ceci : pour
\diagram{\begin{tikzcd}[ampersand replacement=\&]
X\arrow[d,"f"'] \& X'\arrow[l,"q"']\arrow[d,"f'"] \\
Y \& Y'\arrow[l,"p"]
\end{tikzcd}}
cartésien dans $\mathrm{Cat}$, si $f\in\mathrm{Fib}_W$, $p\in W$ alors $q\in W$.

$W(4a')$ Comme a), en supposant $\unsure{f|S}\in\mathrm{Parf}_W$, i.e. $f$ $W_\infty$-parfait (donc une $W_\infty$-fibration) où $W_\infty\subset W$ est la plus petite partie $\unsure{\bar Z|\bar z}$ de $\mathrm{Fl}(\mathrm{Cat})$ satisfaisant $W(1,2\,\mathrm{bis},3)$.

\tikz[baseline=(n.base)]\node[draw,ellipse,inner sep=2pt](n){$W(4b)$}; Soit un diagramme comm.
\diagram{\begin{tikzcd}[ampersand replacement=\&]
X\arrow[rr,"f",no head]\arrow[dr,"p"',no head] \& \& Y\arrow[dl,"q",no head] \\
\& S \&
\end{tikzcd}}
dans $\mathrm{Cat}$, avec $p,q\in\mathrm{Fib}_W$, $f\in W$. Alors $\forall s\in S$, $f_s:X_s\unsure{\relbar|\longrightarrow}Y_s$

$W(4b')$ Comme b), en supposant \deleted{$p$ et $q$ \unsure{fibr}} $X$ et $Y$ fibrés scindés sur $S$, $f$ compatible aux scindages (donc foncteur cartésien).

$W(4b'')$ Comme b), en supposant $p,q$ $W$-\uline{parfaits}

\deleted{$W(4b''')$ Comme b), en supposant $p,q$ $W_\infty$-parfaits}

\tikz[baseline=(n.base)]\node[draw,ellipse,inner sep=2pt](n){$W(4c)$}; Soit $f:X\longrightarrow Y$, avec $f\in\mathrm{Fib}_W\cap W$, alors les fibres de $f$ sont $W$-asphériques

\tikz[baseline=(n.base)]\node[draw,dashed,ellipse,inner sep=2pt](n){$W(4c')$}; Comme c, en supposant $f\in\mathrm{Parf}_W\cap W$. (C'est l'ancien $W(4b'')$, dépouillé des $W_\infty$ parasites)
% archive sheet number: 104
\scanpage{105}
\mspage{54}
J'ai cerclé les notions qui ont été vraiment utiles. Pratiquement, on peut parier que c'est la notion la plus forte, $W(4a)$, qui est la seule importante.
Voici les implications tautologiques (à l'exception de $W(4a')\Longrightarrow W(4b)$, qui est \deleted{\unsure{la principale}} \inserted{\unsure{démontré|démontrée}} dans la dém. de la prop. 4 (pages 6–10).
% slanted
\margin{Cette démonstration utilise une notion plus forte de $W$-fibration que celle du th. 3 p. 51}
\diagram{\begin{tikzcd}[ampersand replacement=\&, row sep=small,column sep=small]
\& a\arrow[d,Rightarrow] \& {} \\
\& a'\arrow[d,Rightarrow] \& {} \\
\& b\arrow[dl,Rightarrow]\arrow[d,Rightarrow]\arrow[dr,Rightarrow] \& {} \\
b' \& c\arrow[d,Rightarrow] \& b'' \\
\& c' \& {}
\end{tikzcd}}
% slanted
\margin{implication douteuse il faut retoucher un peu b), en renforçant la notion de $W$-fibration utilisée pour énoncer b)}

J'ai surtout sué pour montrer que $c'\Rightarrow c$, ou $c\Longrightarrow b$, sans y arriver. L'effort m'a permis de me débarrasser de diverses idées fausses, et de les remplacer par des justes, cf plus bas. Il y a quand même un résultat positif qui surnage :

\uline{Proposition 1} : La condition $b$ (et sans doute même $b''{}^*$, mais je ne vais pas le prouver \unsure{\ldots}) implique \inserted{sans doute} $a$ pour $\bar W$, le saturé fort de $W$. \deleted{\unsure{D}} \inserted{En tous cas,} si $W$ est fortement saturé, les conditions $a,a',b$ (et sûrement aussi $b'$ et $b''$) sont équivalentes.
% slanted
\margin{et en supposant, dans le cas de $b''$, que $W$ satisfait la condition \uline{$b'$} des Limites inductives $W(6\,\mathrm{bis})$. cf. \deleted{le précédent} prop. 3 de ``\unsure{variations du|variations sur le|corollaire de} \unsure{\ldots}'' page 14}

Démonstration de $b\Longrightarrow a$, \deleted{\unsure{pour $W$}} \inserted{quand $W$ est ft saturé.} \deleted{\unsure{Il}} En fait, moyennant b) (et sans $=$ supposer $W$ ft \unsure{saturé|saturée}) nous allons prouver que si on a
\diagram{\begin{tikzcd}[ampersand replacement=\&]
X\arrow[d,"f"'] \& X'\arrow[l,"q"']\arrow[d,"f'"] \\
Y \& Y'\arrow[l,"p"]
\end{tikzcd}}
cartésien dans $\mathrm{Cat}$, avec
% archive sheet number: 105
\scanpage{106}
\mspage{\deleted{\unsure{8}}55}
$f$ une $\overset{\deleted{\unsure{-}}}{W}$-fibration, et $p\in\overset{\deleted{\unsure{-}}}{W}$, alors $\underline{q\in\bar W}$, i.e. $\mathrm{hot}_W(q)$ iso. Pour ceci, nous interprétons la flèche $q$ dans $\mathrm{Hot}_W$ \deleted{\unsure{comme}} de la façon suivante.
Soit
\[
\unsure{\xi|\zeta}=\mathrm{hot}_{\unsure{Y|S},W}(X/Y),\qquad \xi'=p^*(\deleted{\unsure{Y}}\xi),
\]
je dis que $q$ s'identifie à la flèche canonique
\[
H_{\unsure{\bullet|\cdot}}^{\mathrm{HOT}_W}(Y',\xi')\longrightarrow H_{\bullet}^{\mathrm{HOT}_W}(Y,\xi)\quad\text{dans }\substack{\mathrm{Hot}_W\\ \Vert\\ \mathrm{HOT}_W(e)}\tag*{\deleted{\unsure{a}} (*)}
\]
où on note de plus que $\xi\in\mathrm{HOT}_W(\unsure{*|Y})$ est un coeff. \uline{loc. constant} sur $Y$.

Établissons d'abord cette interprétation de $q$. Notons que comme $X\in\underline{\mathrm{Fib}}_W Y$, l'image inverse $\xi'=p^*(\xi)$ n'est autre que $\mathrm{hot}_{Y',W}(X'=X\times_Y Y')$.
D'autre part, si \deleted{$g,g'$ sont} $g:Y\unsure{\relbar|\longrightarrow}e$, $g':Y'\longrightarrow e$ sont les flèches structurales, les deux \unsure{mbs|mbres} de (*) s'interprètent \inserted{\unsure{respt}} comme $g'_!(\xi')$ et $g_!(\xi)$ \deleted{Or le} (les $g'_!,g_!$ au sens du prédérivateur $\mathrm{HOT}_W$).
Or le calcul des $g'_!,g_!$ nous donne que les deux mbs s'identifient resp. à
\[
\mathrm{hot}_W(X'),\qquad\mathrm{hot}_W(X),
\]
et nous admettrons que la flèche \deleted{\unsure{du}}
\[
\mathrm{hot}_W(\unsure{X''|X'})\longrightarrow\mathrm{hot}_W(\unsure{X|x})
\]
n'est autre que $\mathrm{hot}_W(q)$. D'ailleurs, \deleted{\unsure{si}} \unsure{puisque} $X$ est $W$-fibré sur $S$, par définition pratiquement $\xi$ est dans $\mathrm{HOT}_W^{\mathrm{lc}}(S)$. Donc il nous suffira de prouver la
% archive sheet number: 106
\scanpage{107}
\mspage{56}
\uline{Corollaire} Supposons que $W$ \unsure{satisfais|satisfait|satisfasse} à $W(4b)$, et soit $p:S'\longrightarrow S$ dans $W$, $\unsure{\xi|\zeta}\in\mathrm{HOT}_W^{\mathrm{lc}}(S)$, $\xi'=p^*(\xi)$. Alors
\[
H_{\unsure{\bullet|\cdot}}^{\unsure{k|K}}(S',\xi')\overset{\sim}{\longrightarrow}H_{\bullet}(S,\xi)\qquad\text{est un iso. de }\mathrm{Hot}_W.
\]

\uline{Dém} : Il suffit de prouver que pour \unsure{ts} $y\in\mathrm{Ob}\,\mathrm{Hot}_W$, on a une bijection pour les $\mathrm{Hom}(-,y)$. Par les formules d'adjonction sur $S,S'$, cette application entre les $\mathrm{Hom}$ de $\mathrm{Hot}_W$ s'identifie à
\[
\mathrm{Hom}_{\mathrm{HOT}_W(S)}(\xi,y_S)\longrightarrow\mathrm{Hom}_{\mathrm{HOT}_W(S')}^{\unsure{\in|k}}(\xi',y_{S'}),
\]
où $y_S=p^*(y)$, $y_{S'}^{\unsure{s}}=p'{}^*(y)$, donc $y_{S'}=p^*(y_S)$, et la flèche \unsure{(**)} est la flèche sur les $\mathrm{Hom}$ induite par le foncteur
\[
p_{\mathrm{lc}}^*:\mathrm{HOT}_W^{\mathrm{lc}}(S)\longrightarrow\mathrm{HOT}_W^{\mathrm{lc}}(S').
\]
Or moyennant l'axiome $W(4b)$\unsure{,} (que nous n'avons pas utilisé jusqu'à présent), ledit foncteur est une \uline{équivalence de catégories} si $\unsure{p|f}\in W$ (cf \uline{VI} p. 204) : \unsure{qed.|qct.}

De peur de les oublier, je vais maintenant faire une liste des petites questions que j'ai élucidées chemin faisant \unsure{au gribouillis|en gribouillant}.

1°) Si $X,Y$ sont $W$-lisses sur $S$, et $f:X\unsure{\relbar|\longrightarrow}Y$ un $S$ morphisme tel que $f_s:X_s\longrightarrow Y_s$ soit $W$-lisse pour tout $s\in\mathrm{Ob}\,S$, il ne s'ensuit \uline{pas} que $f$ soit $W$-lisse. Mais c'est pourtant vrai quand $X$ et $Y$ sont \unsure{fibrés|fibrées} \unsure{scindés|scindées} sur $S$ (ce
% archive sheet number: 107
\scanpage{108}
\mspage{57}
qui implique déjà leur $W$-lissité sur $S$), et $f$ compatible avec les scindages et a fortiori cartésien. Énoncé dual pour les catégories cofibrées \inserted{scindées} et les foncteurs str. cocartésiens entre elles qui sont $W$-propres fibre par fibre : ils sont $W$-propres.

2°) J'étais intéressé, dans un cas de catégories fibrées scindées \inserted{$X,Y$ sur $S$} et d'un foncteur strictement cartésien $f:X\unsure{\relbar|\longrightarrow}Y$, supposé $W$-parfait fibre par fibre, de conclure que $f$ soit aussi $W$-parfait. \deleted{\unsure{Je savais} \unsure{Comme} \unsure{par 1°)}} (C'\deleted{\unsure{était}}\unsure{est|était} un cas où de plus on sait \deleted{\unsure{déjà}} que
$X,Y$ sont aussi des $W$-fibrés sur $S$, i.e. les foncteurs chgt de base $X_{s'}\longrightarrow X_s$, $Y_{s'}\longrightarrow Y_s$ sont dans $W$. Mais ça n'a pas beaucoup servi \ldots) La question qui restait, vu \uline{1°}, était celle de la $W$-\uline{propreté} de $f$. Mais c'est sans espoir, \inserted{(même pour $Y=S$)} en vertu de \unsure{la|ce} \unsure{brillante|brillant}

\uline{Proposition 2} Supposons que
\diagram{\begin{tikzcd}[ampersand replacement=\&, column sep=small,row sep=tiny]
\& a \\
S=\unsure{o|0}\arrow[ur]\arrow[dr] \& {} \\
\& b
\end{tikzcd}}
, et soit $\deleted{\unsure{X}}Z$ une catégorie sur $S$\deleted{\unsure{.}} \deleted{(prenons $Y=S$)}. Alors $\Phi(Z)$ est propre sur $S$ sss $X=\varnothing$.

Prenons alors dans 2° ci-dessus $X=\Phi(Z)$, $Y=S$, $f:\Phi(Y)\longrightarrow S$ le morphisme structural (donc les $f_s$ sont $W$-parfaits), alors $f$ n'est propre que si $X=\varnothing$. (Pourtant, si on prend $\deleted{\unsure{X}}Z$ $W$-fibré sur $S$, $\Phi(Z)$ l'est aussi, et on a un \unsure{contreexple} à 2°).

Mais on \unsure{s'est|c'est} quand même convaincu du résultat suivant :

\uline{Prop. 3} Soient $X,Y$ cat. fibrées scindées sur $S$, $f:X\longrightarrow Y$ un foncteur str. cartésien, tel que
% archive sheet number: 108
\scanpage{109}
\mspage{58}
pour tout $s\in S$, $f_s:X_s\unsure{\relbar|\longrightarrow}Y_s$ soit $W$-\uline{fibrant}. Alors, \uline{si $W$ satisfait} $W(4a)$, $f$ est $W$-fibrant.
\margin{\unsure{b}}

(J'espérais avoir que $f$ est $W$-fibrant ``à l'œil'', pour prouver que $W(4\unsure{c|C})$ implique $W(4b)$. Mais j'étais obligé de supposer $W(4a)$ \ldots)
% archive sheet number: 109
\scanpage{110}
\mspage{59}
V Sommes amalgamées et carrés $W$-cocartésiens dans \inserted{$\mathrm{Cat}$}

\uline{1}. Généralités sur $U\overset{i}{\lhook\joinrel\relbar}X\overset{j}{\relbar\joinrel\rhook}T$, et les sommes amalgamées $\unsure{X'\amalg_U U'|X'\amalg_U T'}$.

1.1. \inserted{\uline{Données de recollement} entre un ouvert et \deleted{son complémentaire} un fermé :} \deleted{Soit} $X$ dans $\mathrm{Cat}$, $U$ un ouvert de $X$, $T$ le complémentaire fermé.
[ \uline{NB} les ouverts correspondent bien aux ouverts de $\mathrm{Ob}\,X$, topologisé \unsure{pour|par} \unsure{la|sa} relation de \deleted{\unsure{pré}} \inserted{préordre} naturelle $x\leq y\Longleftrightarrow\mathrm{Hom}_X(x,y)\ne\varnothing$, de $=$ pour les fermés.
Les topologies associées aux relations de préordre sont celles où toute intersection d'ouverts est encore ouverte, i.e. toute réunion de fermés est encore fermée. Donc les ouverts d'une telle topologie sont les fermés d'une topologie duale, correspondant à la relation de préordre opposée. La relation de préordre associée à la topologie est
\[
x\leq y\Longleftrightarrow x\in\overline{\{y\}},
\]
la topologie associée au préordre \inserted{sur $I$} est celle pour laquelle les ouverts sont les réunions d'ens de la forme $I_{/x}$ (lesquels forment une base d'ouverts), \inserted{ou encore les ``\uline{cribles}'' ($x\in U,y\leq x\Rightarrow y\in U$),} dont les \unsure{fermés} sont les réunions d'ens $ {}_{x\backslash}I$, c'est-à-dire les cocribles.
Pour \unsure{tt} partie \unsure{$A$} \illegible{1} de $I$, il y a un plus petit ouvert contenant $A$, qui est aussi le plus petit voisinage, c'est la réunion des \deleted{\unsure{$S$}} $I_{/x}$ pour $x\in A$. Et dualement il y a un plus plus grand fermé contenu dans $A$, \deleted{\unsure{d}}dont le complémentaire est le plus petit ouvert contenant $B=I\backslash A$, donc $\bigcup_{x\unsure{\in} B}I_{/x}$. L'adhérence de $A$ \unsure{est} la réunion des adhérences \deleted{de ses} $ {}_{x\backslash}I$ de ses pts. ]
% archive sheet number: 110
\scanpage{111}
\mspage{60}
La situation $(X,U,T)$ se reconstitue, à isom. unique près, par la connaissance de

1°) les deux catégories $U,T$

2°) Le bifoncteur
\[
\left\{\begin{aligned} &H:U^{\unsure{\circ|o}}\times T\unsure{\relbar|\longrightarrow}\mathrm{Ens}\\ &H(u,t)=\mathrm{Hom}_X(u,t)\end{aligned}\right.
\]

De plus, ces données 1°) 2°) ne sont soumises à aucune restriction. On a, plus exactement

\uline{Proposition 1} Un isomorphisme entre la catégorie $\unsure{\mathscr T|\mathcal T}$ dont les objets sont les triples $(X,U,T)$, avec $X\in\mathrm{Ob}\,\mathrm{Cat}$, $U\in\mathrm{Ouv}(X)$, $T\in\mathrm{Ferm}(X)$, $U$ et $T$ complémentaires l'un de l'autre, les \deleted{flèches étant les morphismes} \inserted{morphismes de}
$(X,U,T)$ dans $(X',U',T')$ étant les foncteurs $f:X\longrightarrow X'$ qui induisent $U\longrightarrow U'$, $T\longrightarrow T'$ (\unsure{d'-à|c-à-d|d'à d} \unsure{$U=f^{-1}(U')$}, $T=f^{-1}(T')$), avec la catégorie \inserted{$\mathscr T'$} des triples $(U,T,H)$, de deux catégories $U,T$ et d'un bifoncteur
\[
H:U^\circ\times T\longrightarrow\mathrm{Ens}
\]
(ou encore un préfaisceau $U^\circ\longrightarrow T^\vee$ sur $U$, donc un objet de $\underline{\mathrm{Hom}}(U^\circ,T^\vee)$), un morphisme de $(U,T,H)$ dans $(U',T',H')$ étant un triple
\[
(f_U,f_T,u)\quad\left|\begin{aligned}&\left.\begin{aligned}f_U &: U\longrightarrow U'\\ f_T &: T\longrightarrow T'\end{aligned}\right|\ \text{des foncteurs}\\ &u:H\longrightarrow H'\circ(f_U^\circ\times f_T)\quad\text{i.e.}\\ &\quad H(u,t)\longrightarrow H'(f_U(u),f_T(t)\\ &\quad\text{un hom. fonctoriel (flèche}\\ &\quad\text{dans }(U\times T^\circ)^\wedge\simeq\underline{\mathrm{Hom}}(U^\circ,T^\vee)\end{aligned}\right.
\]
composition des flèches évidente.

Si $U\mathrel{\lhook\joinrel\relbar}X\mathrel{\relbar\joinrel\rhook}T$ est un triple, \deleted{\unsure{d'une}} et
\[
f_{\unsure{U'|U}}:U'\longrightarrow U,\ f_{\unsure{T'|T}}:T'\longrightarrow T
\]
des flèches de $\mathrm{Cat}$, il y a une façon privilégiée de ``recoller'' $U'$ et $T'$, de façon à ce que le l'objet $(X',U',T')$ s'envoie dans $(X,U,T)$ \deleted{à \unsure{l'aide des}} en induisant
% archive sheet number: 111
\scanpage{112}
\mspage{61}
$f_{U'},f_{T'}$ :
\diagram{\begin{tikzcd}[ampersand replacement=\&]
U'\arrow[r,hook,no head,"i'"]\arrow[d,"{f_{U'}}"'] \& X'\arrow[d,"f"] \& T'\arrow[l,hook',no head,"j'"']\arrow[d,"{f_{T'}}"] \\
U\arrow[r,hook,"i"] \& X \& T\arrow[l,hook',"j"']\quad ;
\end{tikzcd}}
elle s'obtient à l'aide du bifoncteur
\[
\left\{\begin{aligned} H'(u',t')&\overset{\mathrm{def}}{=}H(f_{U'}(u'),f_{T'}(t'))\\ \text{i.e.}\qquad H'&= H\circ(f_{U'}^{\unsure{\circ|o}}\times f_{T'}).\end{aligned}\right.
\]
\unsure{Cet|Ce} triple \inserted{i.e. cet objet de $\unsure{\mathscr T|\mathcal T}$} représente le foncteur suivant sur $\mathscr T$ :
\[
(X'',U'',T'')\longmapsto\mathrm{Hom}_{\mathrm{Cat}}(U'',U)\times\mathrm{Hom}_{\mathrm{Cat}}(T'',T),
\]
i.e. on a pour tout objet $(X'',U'',T'')$ de $\mathscr T$ une bijection can., fonctorielle en cet objet
\[
\mathrm{Hom}_{\mathscr T}((X'',U'',T''),(X',U',T'))\simeq\mathrm{Hom}_{\mathrm{Cat}}(U'',U)\times\mathrm{Hom}_{\mathrm{Cat}}(T'',T).\tag*{(*)}
\]
Nul doute qu'on a même
\[
\underline{\mathrm{Hom}}_{\mathscr T}((X'',U'',T''),(X',U',T'))\simeq\underline{\mathrm{Hom}}(U'',\unsure{U|U'})\times\underline{\mathrm{Hom}}(T'',\unsure{T|T'}),\tag*{(**)}
\]
tout aussi fonctoriel, dont (*) se déduit en passant aux ens. d'objets. On dit que le triple $(X',U',T'')$ \deleted{$X'$ est \unsure{dé}} est \deleted{déduit} le triple \uline{induit} par le triple $(X,U,T)$, via les foncteurs $f_{\unsure{U'|U}},f_{T'}$.
\margin{\uline{NB} Si $f_{U'},f_{T'}$ sont pl. fidèles, alors il en est de même de $f:X'\longrightarrow X$. Nous utiliserons cette opération de recollement privilégié surtout quand $U'=U,f_{U'}=\mathrm{id}_U$, ou $T'=T,f_{T'}=\mathrm{id}_T$.}

\deleted{\unsure{Rel}} 1.2. \uline{Relation avec le recollement des topos.}
La situation d'un triple \deleted{d'un}
\[
(\unsure{\mathscr X|\mathcal X},\mathscr U,\mathscr T)
\]
d'un topos, $\mathscr X$, d'un ouvert $U$ de $\mathscr X$, et d'une
% archive sheet number: 112
\scanpage{113}
\mspage{62}
sous-topos fermé complémentaire, est toute similaire.
(\uline{NB} Si $\unsure{\mathscr X|\mathcal X}=\mathrm{Top}(X)$, avec $X$ dans $\mathrm{Cat}$, alors $\mathrm{Ouv}(\mathscr X)\simeq\deleted{\unsure{\mathrm{Ouv}}}\mathrm{Ouv}(X)$, puisque les ouverts correspondent aux sous-faisceaux du faisceau final sur $\mathscr X$, i.e. du préfaisceau final sur $X$. Cette bijection est compatible avec le passage aux complémentaires.)
La théorie des topos dit que les couples $(\mathscr X,\mathscr U,\unsure{\mathscr T|\mathcal T})$ se réconstituent, à \unsure{équiv.|équi.} définie à isom unique près, à l'aide d'un triple
\[
\left\{\begin{aligned}&\mathscr U,\mathscr T&&\text{deux topos}\\ &\varphi:\mathrm{Fais}(\mathscr U)\longrightarrow\mathscr T &&\text{un foncteur exact}\\ &&&\text{à gauche (``accessible''),}\end{aligned}\right.
\]
ce foncteur $\varphi$ étant obtenu, quand on part de
\diagram{\begin{tikzcd}[ampersand replacement=\&]
\mathscr U\arrow[r,"i"] \& \mathscr X \& \mathscr T\arrow[l,"j"'] ,
\end{tikzcd}}
comme
\[
\varphi=j^*i_*.
\]
De plus, ce triple $(\mathscr U,\mathscr T,\varphi)$ peut être choisi arbitrairement, avec la seule restriction (pratiquement anodine, car toujours satisfaite pour tous les foncteurs naturels qu'on rencontre) que $\varphi$ soit (en plus d'être exact à gauche) ``accessible''.

Dans le cas actuel \deleted{\unsure{$\mathscr U$|$\mathscr X$}} $\mathscr X=\mathrm{Top}\,X,\mathscr U=\mathrm{Top}\,U,\mathscr T=\mathrm{Top}(T)$, le foncteur \deleted{\unsure{$\mathrm{Top}$}} $\varphi=j^*i_*$ \deleted{$\mathscr X$ est} non seulement exact à gauche, mais même il commute aux $\varprojlim$ quelconques, puisque il en est ainsi (toujours) de $i_*$, et aussi (puisque \inserted{le} $j$ \inserted{\unsure{toposique}} est induit par le $j:T\longrightarrow X$ catégorique) de $j^*$. En fait on
% archive sheet number: 113
